\chapter{A matched virtual laboratory for calibrated inference}
\label{chap:ci-calibrated-benchmark}

This study asks two separate questions. Does genuine detector calibration
expose a distinction hidden from a specified positive-lag menu? Once the same
finite-data confidence region is available, does nonlinear inference give a
stronger quantitative lower bound than the optimized quadratic calculation?
A third-lag alternative tests whether calibration is the best use of the same
number of additional pair records. Its result is retained even if it outperforms
the calibrated arm. All evidence here comes from simulated trajectories and
readouts, not a physical experiment.

\section{Frozen model, records and comparisons}
The model is the twenty-state ring from Chapter~\ref{chap:ci-nonlinear}, with
$p(1)=0.69$, $p(2)=0.59$, $p(3)=0.485712$, mean zero and conditional detector
second moment $A=1/2$. A record starts at an independently sampled uniform
state. Independent Poisson forward and backward counts give its exact
continuous-time random-walk displacement at the declared duration. The two
binary readings are then drawn independently from their state-dependent
probabilities. Thus the simulator generates hidden paths and detector outputs;
it does not sample agreement bits from the analytic endpoint formula.
Floating-point rates and random-number generation define the numerical
simulation. An independent matrix-exponential calculation checks its endpoint
laws and supplies known-truth coverage diagnostics; it is not supplied to
inference as a measured calibration.

The contract was frozen before production under seed $73024119$. It specifies
$400$ independent datasets in each of four cases: the driven ring; a reversible
ring with the same total jump activity; the driven ring with detector gain
$0.2$; and a reversible ring whose calibration readouts have gain $0.2$ while
its science readouts have gain one. The last case deliberately violates the
fixed-detector contract. It tests whether an incorrect calibration premise can
produce an apparently convincing certificate, not the advertised statistical
coverage under valid premises.

Each calibrated dataset contains $8000$ pairs at lag one, $8000$ at lag two,
and $4000$ same-state calibration pairs. A calibration pair uses two genuinely
independent noise draws conditioned on the same unchanged hidden state.
The mean and second-moment intervals jointly fail with probability at most
$0.025$; the predeclared science contrast $(38/25,-1)$ has one-sided failure
probability $0.025$. Their union therefore gives a $0.05$ total bound.
The nonlinear certificate and the optimized Pinsker certificate use the same
covered halfspace, calibration polygon, counts and path durations. Nonlinear
mixture weights are fixed rational approximations to the quadratic method's
optimal time weights. Its continuous-domain dual check and feasible primal
bracket remain authoritative if numerical score optimization is imperfect.

The alternative arm keeps the same first two settings and spends its $4000$
additional pair records at lag three. It receives the same analytical prior
model knowledge used to choose the calibrated contrast. Before any production
sample, the fixed grid $s=j/100$, $j=1,\ldots,99$ selected the contrast
\[
 a_s=(-s^2,2s,-1)
\]
maximizing its model-centered finite-data Pinsker lower rate at the actual
quotas and durations. The support follows from
$-x(x-s)^2\le0$, giving $-(1-s)^2/2$ for the uncalibrated class.
The selected value is $s=17/50$, with coefficients
$(-289/2500,17/25,-1)$. The inherited $s=2/3$ comparison is also reported.
These are separate predeclared tests at level $0.05$, not an uncorrected
data-selected union test.

\section{Additional access and its cost}
Both arms use $20000$ endpoint pairs and $40000$ readings. The calibrated
arm's declared evolution time is $24000$ model units; the extra-lag arm's is
$36000$. This matches pair/readout counts, not total experimental time.
Ideal same-state calibration has zero intervening evolution in this simulator;
its detector exposure, reset time, two-channel hardware and energy are not
calibrated and cannot be treated as zero physical cost. Conversely, the
extra-lag arm requires a longer evolution window. A hardware conclusion would
require these overheads and access constraints to be measured.

The missing distinction in the original two-lag menu is exact, independent of
Monte Carlo power. A stationary symmetric two-state chain with jump rate
$-\log(9/19)/2$ and detector means $\pm\sqrt{361/450}$ has zero entropy
production and agreement probabilities
\[
 q_k={1+(361/450)(9/19)^k\over2}.
\]
It gives $q_1=0.69$, $q_2=0.59$, exactly matching both joint binary-pair laws
because their means are zero. Its detector second moment is $361/450$, which
differs from the driven model's $1/2$. Genuine same-state calibration can
exclude it. A third positive lag can also distinguish these examples, so this
construction proves a limitation of the stated two-lag access, not of every
uncalibrated experiment.

\section{Held-out findings, including the failed comparison}
The following counts are positive certificates among $400$ independently
generated datasets per case. ``Nonlinear'' means the raw frozen numerical
proposal and exact dual checker, before any fallback to the quadratic result.
\begin{center}
\begin{tabular}{lrrrr}
\hline
Case & Cal. Pinsker & Cal. nonlinear & Extra lag & Legacy extra lag\\
\hline
Driven &350&348&400&392\\
Equilibrium &0&0&0&0\\
Readout gain $0.2$ &0&0&5&5\\
Changed calibration channel &400&400&0&0\\
\hline
\end{tabular}
\end{center}
The last row is outside the fixed-detector contract: those calibrated outputs
are apparent dissipation in a system whose true entropy rate is zero.
They are not successful physical detections. The calibration region contains
the science detector's true moments in none of those $400$ datasets. In the
three valid cases it contains them in all $400$, and the joint science and
calibration event covers in $399/400$ per case. Such Monte Carlo coverage is
a check of these specified simulation laws, not validation of a laboratory's
detector stability or the Markov assumption.

For the driven case, Pinsker's detection rate is $87.5\%$ with pointwise
$95\%$ Clopper--Pearson interval $[83.85\%,90.58\%]$; the raw nonlinear
method gives $87.0\%$ with interval $[83.30\%,90.14\%]$.
The designed extra-lag method detects in all $400$ datasets, with interval
$[99.08\%,100\%]$, and the legacy version in $392/400$.
A zero count in $400$ has upper interval endpoint $0.918\%$; it is not proof
of zero false-positive probability. The comparison demonstrates that this
two-stage calibration protocol is not the best of the tested uses of a fixed
pair budget. The stronger extra-lag arm also consumes $50\%$ more declared
evolution time. Neither comparison determines total energy or hardware cost.

Despite losing two marginal positive decisions through numerical approximation,
the raw nonlinear method raises the all-dataset mean entropy lower bound from
$0.0002355009853$ to $0.0002716097823$, a $15.33\%$ increase in the ratio
of means. This is a same-record quantitative improvement, not improved
detection power in the frozen implementation. All $350$ driven numerical
proposals report solver success, but success is not a proof of zero
optimization error. Their certified relaxation-gap mean is
$1.15\times10^{-7}$, with maximum $2.40\times10^{-6}$.
When the conservative support check yields no positive gap, the implementation
returns no certificate without solving the projection; such skipped cases are
not reported as exact zero KL projections.

The true driven entropy rate is approximately $3.55736$. The mean nonlinear
bound certifies only $0.007635\%$ of that value. The designed extra-lag mean
bound is larger, $0.0010998649243$, but still only about $0.03092\%$ of the
true rate. Thus the study improves quantitative inference while exposing a
substantial remaining loss of information. It does not reconstruct hidden
dissipation.

Because the calibrated quadratic and nonlinear bounds hold on the same
covered event, their maximum is also a valid lower bound without another
failure-probability allocation. A separately labeled audit supplement reports
this practical fallback while preserving every frozen raw result. It replays
the retained counts, stores each computed primal--dual gap, and uses a finer
proposal grid only to diagnose the marginal numerical sign losses. These
post-production diagnostics are not substituted into the frozen comparison.

\section{Reproducibility and limits}
The contract and seed are recorded in
\path{contracts/calibrated_benchmark.json}. The count archive retains every
dataset's science and calibration counts and every method's rate estimate.
The first driven dataset additionally retains hidden endpoints, forward and
backward Poisson counts, elapsed durations and both raw detector readings.
The JSON report gives pointwise $95\%$ Clopper--Pearson Monte Carlo intervals
for detection and coverage, representative exact certificates, all optimized
relaxation gaps, and the simulator's true entropy rate.

A confidence region that covers the observations under an assumed detector
does not validate that detector contract. The gain-mismatch control makes this
distinction explicit. Even a positive, accurately optimized nonlinear lower
bound does not reconstruct the total dissipation of the driven ring. The
quantitative report compares its magnitude with known simulated entropy and
retains the gap. No result here establishes a new universal physical law,
empirical heat dissipation, or a ranking across all sensing instruments.
